\documentclass[fleqn,leqno]{article}

\usepackage[a4paper]{geometry}
\usepackage{parskip}
\usepackage{macros}

\title{Initial Models via a Modality}
\author{Owen Lynch}

\begin{document}

\maketitle

\section{Overview}

In this document, we give a sketch of how to incorporate initial models into geolog.

As suggested by Bob Atkey, we do this via a modality that controls how we are permitted to use query-types in the context when forming theories that are to have initial models.

For this, we follow Shulman's account of modal type theory from (\cite{shulman-2023-semantics}). Shulman's type theory is very general because it is parameterized by an arbitrary ``mode theory.'' Fortunately, as the mode theory that we use is quite simple, we can dispense with a lot of the complexity of the general setup, and use more tightly tailored notation as well.

We will start in \cref{sec:matt} by reviewing the rules of multimodal adjoint type theory, as specialized to our concrete use case. Then in \cref{sec:initial} we will discuss how to use this modality to restrict the formation of initial models. Finally, in \cref{sec:semantics} we will propose a semantics.

\section{Review of multimodal adjoint type theory} \label{sec:matt}

Multimodal adjoint type theory is parameterized by a 2-category. In our case, we are parameterized by the 2-category $\mc{A}$ with the following parts.

\begin{enumerate}
  \item Objects $c$ and $i$ called ``the conjunctive mode'' and ``the inductive mode'' respectively.
  \item A morphism $l \colon c \to i$ called the ``lock'' along with its right adjoint $l^\dagger \colon c \to i$ called the ``unlock.''
  \item We have equations $l \circ l^\dagger = 1_i$ and the counit $\epsilon \colon l \circ l^\dagger \to 1_i$ is the identity 2-morphism, but the unit $\eta \colon 1_c \to l^\dagger \circ l$ is non-trivial.
\end{enumerate}

We then have copies of the type theory at each mode (object). We denote the mode we are in by subscripting the judgment, so we have judgments
%
\begin{mathpar}
  \Gamma \jctx_c \and \Gamma \jctx_i \and
  \Gamma \yields_c A \jtype \and \Gamma \yields_i A \jtype \and
  \Gamma \yields_c a \colon A \and \Gamma \yields_i a \colon A
\end{mathpar}
%
We move between the modes using the 1-morphisms.
%
\begin{mathpar}
  \inferrule{\mu \colon m \to n \and \Gamma \jctx_n}{\Gamma/_\mu \jctx_m} \and
  \inferrule{\Gamma \jctx_m}{\Gamma/_{1_m} \equiv \Gamma \jctx_m} \and \inferrule{\mu \colon n \to p \and \nu \colon m \to p \and \Gamma \jctx_p}{\Gamma/_\mu/_\nu \equiv \Gamma/_{\mu \circ \nu} \jctx_m}
\end{mathpar}
%
So in particular, if $\Gamma \jctx_i$, then $\Gamma/_l/_{l^\dagger} \equiv \Gamma/_{l \circ l^\dagger} \equiv \Gamma/_{1_i} \equiv \Gamma$. This is what motivates the ``lock/unlock'' terminology; if we ``lock'' $\Gamma$, forming $\Gamma/_l$, then applying $-/_{l^\dagger}$ ``unlocks'' $\Gamma/_l$. The other way around does not become the identity, however, because $l^\dagger \circ l$ is not the identity. Intuitively, unlocking followed by locking produces a locked context.

The context formation rule is then modified to permit cross-mode bindings.
%
\begin{mathpar}
  \inferrule{\Gamma \jctx_n \and \mu \colon m \to n \and \Gamma/\mu \yields_m A \jtype}{\Gamma, x \bindvia{\mu} A \jctx_n}
\end{mathpar}
%
We can also incorporate 2-morphisms into substitutions
%
\begin{mathpar}
  \inferrule{\Gamma \yields_n \theta \colon \Delta \and \mu,\nu \colon m \to n \and \alpha \colon \mu \To \nu}{\Gamma/_\nu \yields_m \theta/_\alpha \colon \Delta/_\mu}
\end{mathpar}
Finally, we give the variable rule (note that $\uparrow^\alpha$ is the weakening substitution slashed with $\alpha$).
%
\begin{mathpar}
  \locks(\cdot_m) = 1_m \and
  \locks(\Gamma, x \bindvia{\mu} A) = \locks(\Gamma) \and 
  \locks(\Gamma/_\mu) = \locks(\Gamma) \circ \mu \\
  \inferrule{\alpha \colon \mu \To \locks(\Delta) \colon m \to n}{\Gamma, x \bindvia{\mu} A, \Delta \yields_m A[\uparrow^\alpha]} \and
  \inferrule{\alpha \colon \mu \To \locks(\Delta) \circ \nu \and \beta \colon \nu \To \varsigma}{(\Gamma, x \bindvia{\mu}, \Delta)/_\varsigma \yields x^\alpha[1_{(\Gamma,x\bindvia{\mu},\Delta)}/_\beta] = x^{(\locks(\Delta) \triangleright \beta) \circ \alpha}}
\end{mathpar}
%
Essentially, any variable in context which was previously ``cross-mode bound'' in the mode that we are \emph{currently in} can be used if we give a 2-cell between the 1-morphism that it was bound via and the 1-morphism given by all of the locks between it and the end of the context. As a special case, we have
%
\begin{mathpar}
  (\Gamma,x \bindvia{\mu} A)/_\mu \yields x := x^{1_\mu} \colon A
\end{mathpar}
%
That is, we omit the 2-morphism if it is the identity.

The positive modality for $l$ then allows us to use a conjunctive type in the inductive mode.
%
\begin{mathpar}
  \inferrule{\Gamma/_l \yields_c A \jtype}{\Gamma \yields_i \lozenge A \jtype} \\
  \inferrule{\Gamma/_l \yields_c a \colon A \jtype}{\Gamma \yields_i [a]_\lozenge \colon \lozenge A} \and
  \inferrule{\Gamma/_{l^\dagger} \yields_i d \colon \lozenge A \and \Gamma, y \bindvia{l^\dagger} \lozenge A \yields_c B \jtype \and \Gamma, x \bindvia{l^\dagger \circ l} A \yields_c b \colon B[y \subst [x]_\lozenge]}{\Gamma \yields_c (\llet [x]_\lozenge \ot d \inn b) \colon B[y \subst d]} \\
  (\llet [x]_\lozenge \ot [a]_\lozenge \inn b) \equiv b[x \subst a]
\end{mathpar}
Crucially, because there is no morphism $l^\dagger \circ l \to 1_c$, we cannot use a variable $x \bindvia{l^\dagger \circ l} A$ in the same way that we would use a variable $x \colon A$. We think of a variable $x \bindvia{l^\dagger \circ l} A$ as an ``inductive model'' of $A$, and a variable $x \colon A$ as a ``conjunctive model.''

\section{Initial models and modal pi types} \label{sec:initial}

Now that we have laboriously set up the modal type theory, the rule for initial models is very simple.
%
\begin{mathpar}
  \inferrule{\Gamma/_l \yields_c A \jtype}{\Gamma \yields_i \init A \colon \lozenge A}
\end{mathpar}
%
This gives both the formation and introduction rules for inductive types (we will see exactly how in a bit). The elimination and computation rules need a lot of machinery to set up, so we will not cover them just yet.

One of the points of this whole machinery was to do initial models \emph{in context}. In order to do this, we still need to use the variables in $\Gamma/_l$ \emph{somehow} in forming $A$. The way this works is that we generalize the pi theories in geott to \emph{modal pi theories}. Assume that we have ``types at level'' judgments $A \jtype_{\ms{query}}$ and $A \jtype_{\ms{theory}}$, as in geott.
%
\begin{mathpar}
  \inferrule{\Gamma/_{l^\dagger} \yields_i A \jtype_{\ms{query}} \and \Gamma, x \bindvia{l^\dagger} A \yields_c B \jtype_{\ms{theory}}}{\Gamma \yields_c (x \bindvia{l^\dagger} A) \to B \jtype_{\ms{theory}}} \\
  \inferrule{\Gamma, x \bindvia{l^\dagger} A \yields_c b \colon B}{\Gamma \yields (\lambda x.b) \colon (x \bindvia{l^\dagger} A) \to B} \and
  \inferrule{\Gamma \yields_c f \colon (x \bindvia{l^\dagger} A) \to B \and \Gamma/_{l^\dagger} \yields_i a \colon A}{\Gamma \yields_c f \, a \colon B[x \subst a]} \\
  (\lambda x.b)\,a \equiv b[x \subst a] \and f \equiv \lambda x. (f\, x)
\end{mathpar}

This allows us to use variables behind the lock in the domain of pi theories, but not in the codomain, which is precisely the restriction we needed for initial models.

\begin{example}
  Sum types are given by the following initial model.
  \begin{mathpar}
    \inferrule{(A \colon \ms{Query}, B \colon \ms{Query})/_{l^\dagger}/_l \yields_c [t \colon \ms{Query}, l \colon A \to t, r \colon B \to t] \jtype}{(A \colon \ms{Query}, B \colon \ms{Query})/_{l^\dagger} \yields_i \init \, [t \colon \ms{Query}, l \colon A \to t, r \colon B \to t] \colon \lozenge [t \colon \ms{Query}, l \colon A \to t, r \colon B \to t]}
  \end{mathpar}
  In order to derive this, we must have rules sufficient to derive $(A \colon \ms{Query})/_{l^\dagger} \yields_i A \jtype_{\ms{query}}$. It is not immediately obvious what sequence of derivations yields this, but it should at least be intuitive that this is permitted.
\end{example}

One thing to note is that ideally $\lozenge A$ would compute through on the various theory formers. E.g., $\lozenge \ms{Query}_c \equiv \ms{Query}_i$, $\lozenge (\Sigma\,A\,B) \equiv \Sigma (\lozenge A) (\lozenge B)$, $\lozenge ((x \bindvia{l^\dagger} A) \to B) \equiv (x \colon A) \to \lozenge B$, etc. Thus, when we get an initial model it is in a form which we can immediately use, via the relevant eliminators.

Another option here would be that $\lozenge (\Sigma\,A\,B)$ \emph{behaves like} $\Sigma\,(\lozenge A)\,(\lozenge B)$, in the sense of the kinetic/potential distinction. This is perhaps ideal when $A$ is itself nominal, because $\lozenge$ can't ``compute through'' here, but $\lozenge$ can act as a ``modifier'' on the behavior of $A$.

Tentatively, the notation for this could look like \verb|ind A|. The intuition for this is that if \verb|a : ind A| then \verb|a| is an inductive model of \verb|A|. So, for instance, if \verb|Either A B| is the initial model of \verb|Either/theory A B|, where
\begin{verbatim}
theory Either/theory (A : Query) (B : Query) := sig
  t : Query
  left : A -> t
  right : B -> t
end
\end{verbatim}
then we would have \verb|Either A B .t : ind Query|.

\section{Semantics} \label{sec:semantics}

This section relies on the assumption that the reader has read the lecture notes on the semantics of geolog.

For the semantics, let $\ms{Conj}$ and $\ms{Ind}$ be two 2-categories, with a 2-adjunction
% https://q.uiver.app/#q=WzAsMixbMCwwLCJcXG1ze0luZH0iXSxbMCwyLCJcXG1ze0Nvbmp9Il0sWzEsMCwiRiIsMCx7ImN1cnZlIjotMn1dLFswLDEsIlUiLDAseyJjdXJ2ZSI6LTJ9XSxbMiwzLCIiLDAseyJsZXZlbCI6MSwic3R5bGUiOnsibmFtZSI6ImFkanVuY3Rpb24ifX1dXQ==
\[\begin{tikzcd}
	{\ms{Ind}} \\
	\\
	{\ms{Conj}}
	\arrow[""{name=0, anchor=center, inner sep=0}, "U", curve={height=-12pt}, from=1-1, to=3-1]
	\arrow[""{name=1, anchor=center, inner sep=0}, "F", curve={height=-12pt}, from=3-1, to=1-1]
	\arrow["\dashv"{anchor=center}, draw=none, from=1, to=0]
\end{tikzcd}\]
such that the counit $\epsilon \colon FU \to 1_{\ms{Ind}}$ is a 2-natural equivalence.

Tentatively, we might have $\ms{Ind}$ be the 2-category of locally small finite limit categories, and $\ms{Conj}$ be the 2-category of arithmetic universes, where $F$ is the \emph{conservative completion} that completes a finite limit category to an arithmetic universes while preserving those finite coproducts/quotients/list objects that may already exist. Another option would be to let $\ms{Ind}$ be the opposite of the category of Grothendieck topoi and geometric morphisms, with $F$ the conservative cocompletion functor.

We will call objects of $\ms{Conj}$ \defcase{conjunctive domains} and objects of $\ms{Ind}$ \defcase{inductive domains}. For a conjunctive domain $\mc{C}$, we call $F(\mc{C})$ its inductive completion (and remember that this is a \emph{conservative} completion).

We then have an induced 2-adjunction in the other direction between the categories of \emph{prestacks} (contravariant functors to $\ms{Cat}$). (Recall from Johnstone that a geometric theory is a prestack on topoi generated by a finite sequence of PIE limits).
% https://q.uiver.app/#q=WzAsMixbMCwwLCJcXG1ze1BzdH0oXFxtc3tJbmR9KSJdLFswLDIsIlxcbXN7UHN0fShcXG1ze0Nvbmp9KSJdLFswLDEsIkZeXFxhc3QiLDIseyJjdXJ2ZSI6Mn1dLFsxLDAsIlVeXFxhc3QiLDIseyJjdXJ2ZSI6Mn1dLFszLDIsIiIsMix7ImxldmVsIjoxLCJzdHlsZSI6eyJuYW1lIjoiYWRqdW5jdGlvbiJ9fV1d
\[\begin{tikzcd}
	{\ms{Pst}(\ms{Ind})} \\
	\\
	{\ms{Pst}(\ms{Conj})}
	\arrow[""{name=0, anchor=center, inner sep=0}, "{F^\ast}"', curve={height=12pt}, from=1-1, to=3-1]
	\arrow[""{name=1, anchor=center, inner sep=0}, "{U^\ast}"', curve={height=12pt}, from=3-1, to=1-1]
	\arrow["\dashv"{anchor=center, rotate=-180}, draw=none, from=1, to=0]
\end{tikzcd}\]
This describes the functor from our mode theory $\mc{A}$ to $\ms{Cat}$. More specifically:
\begin{itemize}
  \item Contexts in mode $i$ are prestacks on $\ms{Ind}$
  \item Contexts in mode $c$ are prestacks on $\ms{Conj}$
  \item The lock $\Gamma \mapsto \Gamma/_l$ is interpreted by $F^\ast$
  \item Unlocking $\Gamma \mapsto \Gamma/_{l^\dagger}$ is interpreted by $U^\ast$
\end{itemize}

Let's use this to work out what the interpretation of the context $(A \colon \ms{Query}, B \colon \ms{Query})/_l$ is. First of all $A \colon \ms{Query}, B \colon \ms{Query}$ is the context which sends an inductive domain $\mc{I}$ to the category of pairs of objects in $\mc{I}$, e.g. $\mc{I} \times \mc{I}$.

Then $(A \colon \ms{Query}, B \colon \ms{Query})/_l$ is the context which sends a conjunctive domain $\mc{C}$ to the category of pairs of objects in $F(\mc{C})$, e.g. $F(\mc{C}) \times F(\mc{C})$. This is why we can't have $(A \colon \ms{Query}, B \colon \ms{Query})/_l \yields_c A \colon \ms{Query}$; we are trying to produce an object of $\mc{C}$ but we only have access to objects of $F(\mc{C})$.

It is currently only a conjecture that the rule for initial models does in fact have semantics in this setup.

\bibliographystyle{plainnat}
\bibliography{geolog.bib}

\end{document}

% Local Variables:
% TeX-engine: luatex
% End:
